% Part: computability
% Chapter: computability-theory
% Section: applications-fixed-point

\documentclass[../../../include/open-logic-section]{subfiles}

\begin{document}

\olfileid{cmp}{thy}{apf}
\olsection{مقرر نقطے دے قضیے دی ورتوں}

مقرر نقطے دے قضیے نال اسی
جزوی قابلِ حساب فنکشناں نوں اوہناں
دے اپنے اشارییاں دے حوالے نال
تعریف کر سکدے آں۔ مثلاً،
اک اشاریہ $e$ لبھ
سکدا اے ایداں کہ
ہر $y$ لئی:
\[
\cfind{e}(y) = e + y.
\]
ہور مثال وچ،
مقرر نقطے دے قضیے
دا ثبوت ورت کے Java
یا C++ وچ ایداں
دا پروگرام بنایا
جا سکدا اے جیڑا
اپنا متن آپ چھاپے۔

یاد رکھو: جے ہر
$e$ لئی $W_e$
نوں $\cfind{e}$
دا دائرۂ تعریف
منئیے، تاں سلسلہ
$W_0$، $W_1$،
$W_2$،~\dots
قابلِ حساب گنتی
وار فہرست بنن جوگ
سیٹاں دی فہرست
دیندا اے۔ ایہناں
وچوں کجھ سیٹ
قابلِ حساب نیں۔
سوال ایہ اے:
کی کوئی الگورتھم
اک قدر $x$
انپٹ لے کے،
جے $W_x$ قابلِ
حساب ہووے،
اوس دے خصوصیتی
فنکشن دا اشاریہ
واپس دے سکدا اے؟
جواب «نہیں» اے؛
ایداں دا کوئی
الگورتھم نہیں:

\begin{thm}
کوئی جزوی قابلِ حساب
فنکشن $f$ نہیں جیہدی ایہ خاصیت ہووے:
جدوں وی $W_e$
قابلِ حساب ہووے،
تاں $f(e)$
تعریف شدہ ہووے
تے $\cfind{f(e)}$
اوس دا خصوصیتی
فنکشن ہووے۔
\end{thm}

\begin{proof}
% PNB-CMP-025: proof must cover partial computable f and undefined-index cases.
من لو $f$ کوئی وی
جزوی قابلِ حساب
فنکشن اے؛ اسی اک
$e$ بناواں گے
ایداں کہ $W_e$
قابلِ حساب ہووے،
پر $\cfind{f(e)}$
اوس دا خصوصیتی
فنکشن نہ ہووے،
یا مطلوبہ اشاریہ
ہی نہ ملے۔
مقرر نقطے دا
قضیہ ورت کے
اک اشاریہ $e$
لبھ سکدا اے
ایداں کہ:
\[
\cfind{e}(y) \simeq
\begin{cases}
0 & \text{جے $y=0$ تے $\cfind{f(e)}(0) \fdefined = 0$} \\
\text{غیر تعریف شدہ} & \text{ہور صورت وچ۔}
\end{cases}
\]
یعنی $e$، مقرر
نقطے دا قضیہ
ہِٹھاں دتے
فنکشن اُتے
لا کے ملدا
اے:
\[
g(x,y) \simeq
\begin{cases}
0 & \text{جے $y=0$ تے $\cfind{f(x)}(0) \fdefined = 0$} \\
\text{غیر تعریف شدہ} & \text{ہور صورت وچ۔}
\end{cases}
\]
غیر رسمی طور اُتے ایویں ویکھ سکدے آں
کہ $g$ جزوی قابلِ حساب اے:
انپٹ $x$ تے
$y$ اُتے الگورتھم
پہلے جانچدا اے
کہ $y$،~$0$
دے برابر اے
یا نہیں۔ جے
اے، تاں اوہ
$f(x)$ حساب
کردا اے،
فیر ہمہ گیر مشین
نال $\cfind{f(x)}(0)$
حساب کردا اے۔
جے ایہ آخری
حساب رُک کے
~$0$ واپس
دیوے، تاں
الگورتھم~$0$
دے کے رُکدا
اے؛ ہور صورت
وچ نہیں رُکدا۔
جے اشاریہ کڈھن والا پہلا حساب
ہی نہ رُکے،
تاں وی پورا
الگورتھم نہیں
رُکدا۔

ہن دھیان دیو:
جے $\cfind{f(e)}(0)$
تعریف شدہ اے
تے $0$ دے
برابر اے، تاں
$\cfind{e}(y)$
عین اُدوں
تعریف شدہ
اے جدوں $y$،
$0$ دے
برابر ہووے؛
اس لئی
$W_e = \{ 0 \}$۔
جے $\cfind{f(e)}(0)$
غیر تعریف شدہ
اے، یا تعریف
شدہ اے پر
$0$ دے برابر
نہیں، تاں
$W_e = \emptyset$۔
جے ایہ اشاریہ موجود ہووے، دونوں
صورتاں وچ $\cfind{f(e)}$،
$W_e$ دا
خصوصیتی فنکشن
نہیں: انپٹ $0$
اُتے صفر آوے
تاں جواب اک
ہونا چاہیدا؛
غیر صفر آوے
تاں جواب صفر
ہونا چاہیدا۔
جے اشاریہ
دین والا حساب
ہی تعریف شدہ
نہ ہووے، تاں
مطلوبہ اشاریہ
نہیں ملدا؛
تے جے اوہ
اشاریہ ملے
پر اوس دا
فنکشن صفر
اُتے غیر تعریف
شدہ ہووے،
تاں اوہ کُل
خصوصیتی فنکشن
نہیں۔
\end{proof}

\end{document}
